\documentclass[10pt,a4paper]{amsart}
\usepackage[margin=19mm]{geometry}
\usepackage{amsmath,amssymb,mathtools}
\usepackage[scr=rsfs,cal=euler]{mathalfa}
\usepackage{xfrac}
\usepackage[unicode,hidelinks,bookmarks=false]{hyperref}
\newcommand{\ie}{i.e.\ }
\newcommand{\cf}{cf.\ }
\newcommand{\resp}{respectively\ }
\newcommand{\Subsection}{\subsection}
\providecommand{\nobreakdash}{\nobreak\hbox{-}}
\setlength{\emergencystretch}{2em}
\multlinegap0pt
\usepackage{amssymb}
\usepackage{mathtools}
\usepackage{tikz}
\usepackage{xcolor}
\usepackage{algorithm}\usepackage{algpseudocode}
\newtheorem{proposition}{Proposition}
\usepackage{cleveref}
\crefname{proposition}{Proposition}{Propositions}
\title{Direct Sampling of Confined Polygons in Linear Time}
\author{Clayton Shonkwiler, Kandin Theis}
\date{Source: arXiv:2501.04885v3 (CC BY 4.0)}
\begin{document}
\maketitle

\noindent\emph{Source and licence.} Direct Sampling of Confined Polygons in Linear Time. Version arXiv:2501.04885v3 (CC BY 4.0), distributed by arXiv under the Creative Commons Attribution 4.0 International licence, http://creativecommons.org/licenses/by/4.0/, as recorded in the arXiv licence metadata for this exact version and shown on the abstract page https://arxiv.org/abs/2501.04885v3. Full licence notice: \texttt{licenses/arxiv-2501.04885-CC-BY-4.0.txt}.\par\smallskip

\noindent\textit{Editorial scope.} Counted unit: the article's proposition prop:phi (source0.tex lines 1267--1277). The article's complete original proof of it is delivered with it (source0.tex lines 1324--1373, one \texttt{proof} environment of 50 lines, which the article places away from the statement). No mathematics is written, repaired or re-derived: the statement and the proof are the article's own lines, in the article's own notation, and no retained line is substituted or re-flowed.  No further source-local context is needed: every cross-reference of every retained line resolves inside the two delivered spans.  Nothing was omitted from any retained span, so the package carries no omission record.  No retained line contains a citation, so the wrapper carries no bibliography and no bibliography entry.  No retained line contains a figure environment, an image-inclusion command or drawing code, so no image is delivered and the package declares no asset dependency.  Editorial formatting only, in addition: an AMS \texttt{amsart} preamble that restates the article's own declarations and macros used by the retained lines, namely the environments \texttt{proposition} on separate counters, together with the standard packages those lines need.  The retained environments print in this document as Proposition 1 in order of appearance, whereas the article prints the same items with the numbering of its own full document.  The complete native source of the article as distributed in the arXiv e-print for this exact version is bundled unchanged as \texttt{sources/171368/source0.tex}.
	\begin{proposition}\label{prop:phi}
		The turning angle $\phi_i$ can be expressed in terms of action-angle coordinates as
		\begin{equation}\label{eq:phi}
			\phi_i = \arccos\left(\frac{d_{i-2}^2+d_i^2}{2} - 1 - \frac{(d_{i-2}^2 + d_{i-1}^2 - 1)(d_{i-1}^2 + d_{i}^2 - 1) - A \cos \theta_{i-1}}{4 d_{i-1}^2}\right),
		\end{equation}
		where 
		\begin{multline*}
			A = \sqrt{(d_{i-2}+d_{i-1}-1)(d_{i-2}-d_{i-1}+1)(-d_{i-2}+d_{i-1}+1)(d_{i-2}+d_{i-1}+1)}  \\
			\times \sqrt{(d_{i-1}+d_{i}-1)(d_{i-1}-d_{i}+1)(-d_{i-1}+d_{i}+1)(d_{i-1}+d_{i}+1)}.
		\end{multline*}
	\end{proposition}

	\begin{proof}[{Proof of \Cref{prop:phi}}]
	Since the edges $e_i = v_{i+1}-v_i$ and $e_{i+1} = v_{i+2}-v_{i+1}$ have unit length, $\cos \phi_i$ is simply their dot product:
	\begin{equation}\label{eq:cosphi1}
		\cos\phi_i = (v_{i+1}-v_i) \cdot (v_{i+2}-v_{i+1}) = v_{i+1} \cdot v_{i+2} - v_i \cdot v_{i+2} - \|v_{i+1}\|^2 + v_i \cdot v_{i+1}.
	\end{equation}
	Since the root vertex $v_1$ is at the origin, 
	\begin{equation}\label{eq:norm vi+1}
		\|v_{i+1}\|^2 = d_{i-1}^2
	\end{equation} 
	and 
	\[
		1 = \|e_i\|^2 = (v_{i+1} - v_i) \cdot (v_{i+1} - v_i) = \|v_{i+1}\|^2 - 2 v_i \cdot v_{i+1} + \|v_i\|^2 = d_{i-1}^2 - 2 v_i \cdot v_{i+1} + d_{i-2}^2,
	\]
	so 
	\begin{equation}\label{eq:consecutive dots}
		v_i \cdot v_{i+1} = \frac{d_{i-1}^2 + d_{i-2}^2 - 1}{2} \quad \text{and similarly} \quad v_{i+1} \cdot v_{i+2} = \frac{d_{i}^2 + d_{i-1}^2 - 1}{2}.
	\end{equation}
	Plugging \eqref{eq:norm vi+1} and \eqref{eq:consecutive dots} into \eqref{eq:cosphi1} yields
	\begin{equation}\label{eq:cosphi}
		\cos \phi_i = \frac{d_{i-2}^2 + d_i^2}{2} - 1 - v_i \cdot v_{i+2}.
	\end{equation}
	
	On the other hand, the dihedral angle $\theta_{i-1}$ is the angle between the unit normal vectors $\frac{v_i \times v_{i+1}}{\|v_i \times v_{i+1}\|}$ and $\frac{v_{i+1} \times v_{i+2}}{\|v_{i+1} \times v_{i+2}\|}$ to the two triangles, so
	\begin{equation}\label{eq:costheta1}
		\cos\theta_{i-1} = \frac{(v_i \times v_{i+1})\cdot(v_{i+1} \times v_{i+2})}{\|v_i \times v_{i+1}\|\|v_{i+1} \times v_{i+2}\|}.
	\end{equation}
	Using the Binet–Cauchy identity, the numerator 
	\begin{multline}\label{eq:quadruple product}
		(v_i \times v_{i+1})\cdot(v_{i+1} \times v_{i+2}) = (v_i \cdot v_{i+1})(v_{i+1} \cdot v_{i+2}) - (v_i \cdot v_{i+2})(v_{i+1} \cdot v_{i+1}) \\
		= \frac{1}{4}(d_{i-2}^2 + d_{i-1}^2 - 1)(d_{i-1}^2 + d_{i}^2 - 1) - d_{i-1}^2 v_i \cdot v_{i+2}
	\end{multline}
	where we've used \eqref{eq:norm vi+1} and \eqref{eq:consecutive dots} to express three of the dot products in terms of the $d$'s.
	
	In the denominator of~\eqref{eq:costheta1}, $\|v_i \times v_{i+1}\|$ is twice the area of the triangle with side lengths $d_{i-2}, d_{i-1}$, and~$1$. We can compute this area using Heron's formula, yielding
	\[
		\|v_i \times v_{i+1}\| = \frac{1}{2} \sqrt{(d_{i-2}+d_{i-1}-1)(d_{i-2}-d_{i-1}+1)(-d_{i-2}+d_{i-1}+1)(d_{i-2}+d_{i-1}+1)}.
	\]
	Substituting this, the analogous formula for $\|v_{i+1}\times v_{i+2}\|$, and \eqref{eq:quadruple product} into \eqref{eq:costheta1} yields
	\begin{equation}\label{eq:costheta}
		 \cos\theta_{i-1}=\frac{(d_{i-2}^2 + d_{i-1}^2 - 1)(d_{i-1}^2 + d_{i}^2 - 1) - 4d_{i-1}^2 (v_i \cdot v_{i+2})}{A}.
	\end{equation}
	
	
	
	Solving~\eqref{eq:costheta} for $v_i \cdot v_{i+2}$ and plugging into~\eqref{eq:cosphi} gives an expression for $\cos \phi_i$:
	\[
		\cos\phi_i = \frac{d_{i-2}^2+d_i^2}{2} - 1 - \frac{(d_{i-2}^2 + d_{i-1}^2 - 1)(d_{i-1}^2 + d_{i}^2 - 1) - A \cos \theta_{i-1}}{4 d_{i-1}^2}.
	\]
	Taking the inverse cosine of both sides completes the proof.
	\end{proof}

\end{document}
